\chapter{Introduction}
\label{ch:introduction}

\draftnote{this chapter is a placeholder skeleton. Expand with a precise
statement of the research questions before submission.}

Ground-penetrating radar (GPR) images subsurface structure by emitting a
short electromagnetic pulse and recording its reflections from dielectric
contrasts in the ground. The achievable image resolution is fundamentally
limited by the wavelength of the probing pulse: two reflectors closer than
roughly half a wavelength apart, or a single reflector that moves by a small
fraction of a wavelength between two surveys, cannot be distinguished from
the migrated \emph{amplitude} image alone (\cref{ch:resolution}). This is a
practical obstacle for time-lapse monitoring applications --- tracking
millimetre-scale ground movement, or the advance of a fluid front through a
sub-wavelength fracture --- where the displacement or material change of
interest is, by construction, far smaller than the wavelength of the radar
pulse used to image it.

This thesis develops and validates a \emph{phase}-based alternative: rather
than reading displacement off the migrated amplitude image, a baseline and a
monitor survey are migrated and compared in the two-dimensional Fourier
domain, where the Fourier shift theorem turns any sub-wavelength translation
into a fully resolvable linear phase ramp (\cref{ch:theory}). A weighted
least-squares fit of this phase plane recovers the displacement with
sub-millimetre precision, and --- critically --- separates a purely
geometric shift from a phase rotation caused by a change in the dielectric
properties of the target itself, which is the signature of, for example, a
fracture filling with water.

\section*{Unifying Hypothesis}

The series of experiments in this thesis tests a single unifying hypothesis,
which no individual experiment can test on its own:

\begin{quote}
\emph{A weighted least-squares fit of the 2D cross-spectrum phase plane
resolves sub-wavelength displacement and material change in GPR time-lapse
surveys, across migration algorithms and under realistic noise and geometry
conditions, more reliably than amplitude-based methods.}
\end{quote}

\Cref{ch:resolution} establishes the amplitude floor this hypothesis claims
to beat. \Cref{ch:tl-horizontal,ch:tl-vertical} confirm that floor persists
under simple amplitude differencing of a genuinely displaced scatterer.
\Cref{ch:phase-plane} tests whether the phase-plane estimator itself recovers
displacement below that floor, across migration algorithms and under noise.
\Cref{ch:fluidflow} tests whether the same estimator still decouples
geometric movement from material change in a more field-realistic, spatially
distributed scenario. \Cref{ch:discussion} returns to this hypothesis in
light of all five experiments together.

\section*{Outline}

\Cref{ch:litreview} reviews the literature this thesis builds on.
\Cref{ch:theory} derives the migration and phase-plane theory used
throughout the thesis. \Cref{ch:methodology} describes the shared simulation
and processing pipeline. \Cref{ch:resolution} establishes the amplitude-based
resolution limit that motivates the phase-based approach.
\Cref{ch:tl-horizontal,ch:tl-vertical} apply simple amplitude differencing to
lateral and vertical sub-wavelength displacements of a point scatterer.
\Cref{ch:phase-plane} validates the 2D phase-plane framework itself, and
\Cref{ch:fluidflow} applies the full pipeline to a more realistic, spatially
distributed fluid-front scenario. \Cref{ch:discussion} closes with a
discussion of the results in light of the unifying hypothesis above, and the
\nameref{ch:summary} restates the thesis's main conclusions.
